\section{Как длина траектории меняет компромисс}
\label{sec:horizon}

За совокупным результатом скрываются различия между задачами с короткими и
длинными траекториями. Используем когорты из раздела~\ref{sec:setup}:
короткую ($T\leq8$), среднюю ($8<T\leq20$) и длинную ($T>20$), где $T$ ---
медианная длина траектории полной истории для данной задачи. Состав каждой
когорты одинаков для всех способов организации памяти. Собственная длина
траектории метода является результатом, а не основанием для его отнесения
к когорте.

\begin{table}[htbp]
\centering\small
\caption{Когорты длины траектории, заданные контролем: качество
и стоимость. Учитываются все включённые исходы. Отношение стоимости ---
медиана отношений к полной истории в той же задаче и том же повторе.}
\label{tab:horizon-main}
\resizebox{\textwidth}{!}{\input{../paper/tables/new-horizon}}
\end{table}

Состояние без истории уступает контролю по успешности во всех когортах.
Наибольший разрыв возникает в длинной: \HorizonLongModelStateRate\% против
\HorizonLongModelControlRate\%. Окно $k=6$ возвращает успешность до
\HorizonLongModelHistorySixRate\% (табл.~\ref{tab:horizon-main}). Отношение
стоимости к контролю составляет \HorizonLongModelHistorySixAllCostRatio{}
против \HorizonLongModelStateAllCostRatio{} у состояния без истории.
Стоимость на наблюдаемый успех равна соответственно
\$\HorizonLongModelHistorySixCostPerSuccess{} и
\$\HorizonLongModelStateCostPerSuccess{}; у контроля ---
\$\HorizonLongModelControlCostPerSuccess{}.

\begin{table}[htbp]
\centering\small
\caption{Те же когорты: медианы контекста и токенов. Cache hit ---
медиана доли $K/(U+K)$ по выполнениям; столбцы токенов --- отдельные
медианы.}
\label{tab:horizon-cache}
\resizebox{\textwidth}{!}{\input{../paper/tables/new-horizon-resources}}
\end{table}

На длинных траекториях усиливается и преимущество полной истории по
использованию кеша. Медианная доля кешированных входных токенов возрастает
с \HorizonShortModelControlCacheHit\% в короткой когорте до
\HorizonLongModelControlCacheHit\% в длинной. У состояния без истории и
при $k=6$ в длинной когорте она составляет
\HorizonLongModelStateCacheHit\% и \HorizonLongModelHistorySixCacheHit\%.
Сам объём кешированных токенов этого различия не показывает: доля зависит
и от некешированного входа.

Когорты показывают, где окно истории полезно, но не задают универсального
порога переключения памяти. У задач короткой когорты встречаются отдельные
длинные повторы; поэтому эффективность на всей нагрузке и медианная
относительная стоимость могут давать разное ранжирование. Траектории
зависят от политики агента. Без контролируемого изменения горизонта задачи
наблюдаемый результат не устанавливает причинной точки перехода.
